% GENERATED FILE. Edit canonical lessons, then run npm run generate:books.
% canonical-source-hash: fnv1a64:2783424cf83088d7
% canonical-lessons: LA-C297-incipio, LA-C297-coepi, LA-C297-finio, LA-C297-perficio, LA-C297-intermitto

\chapter{Incipere et fīnīre}
\label{ch:la-incipere-et-finire}
\hlchaptermodality{297}

\begin{chapteropening}
\textbf{By the end of this chapter:} \emph{I can say that an action began, was finished, was completed or was broken off.}

Complete the last lesson of chapter 297: I can say that an action began, was finished, was completed or was broken off.
\end{chapteropening}

\section[incipiō, incipere]{incipiō, incipere — to begin}

\label{lesson:LA-C297-incipio}



\begin{quote}
Before the new one: say the Latin for a gap, then the Latin for a moment.
\end{quote}

\subsection*{You'll want to know: incipiō, incipere}
\textbf{incipiō, incipere} — \textquotedblleft{}to begin\textquotedblright{}.

\begin{quote}
Mīlitēs mūrum aedificāre incēpērunt.
\end{quote}

\begin{quote}
\emph{The soldiers began to build a wall.}
\end{quote}

Perfect \emph{incēpī}. With an infinitive: \emph{incipit loquī}, he begins to speak. English \emph{incipient}.

\begin{cousinweb}
\textbf{incipiō} is built on \emph{capiō}, \textquotedblleft{}to take\textquotedblright{}, with \emph{in-} (Lewis \& Short): to lay hold of a task, to take it in hand. English \emph{incipient} comes from it.
\end{cousinweb}

\begin{grammarlens}[title={What you've built}]
One more word for telling, explaining and describing at B1.
\end{grammarlens}

\subsection*{Your turn}
\begin{itemize}
\raggedright
  \item \emph{Say it:} \emph{incipiō, incipere}
  \item \emph{Say it:} \emph{incipiō, incipere}, once more
  \item \emph{Say it:} a sentence of your own with \emph{incipiō, incipere}
  \item \emph{Recall:} say the Latin for a gap, then the Latin for a moment, then say \emph{incipiō, incipere} again
\end{itemize}

\begin{tcolorbox}[breakable,colback=teal!4,colframe=teal!35!black,title={Before you move on}]
What does incipiō, incipere mean? (\textquotedblleft{}To begin\textquotedblright{}.) Say it once more.
\end{tcolorbox}
\section[coepī, coepisse]{coepī, coepisse — I began, I have begun}

\label{lesson:LA-C297-coepi}



\begin{quote}
Before the new one: say the Latin for a moment, then the Latin for to begin.
\end{quote}

\subsection*{You'll want to know: coepī, coepisse}
\textbf{coepī, coepisse} — \textquotedblleft{}I began, I have begun\textquotedblright{}.

\begin{quote}
Ubi nox vēnit, puellae canere coepērunt.
\end{quote}

\begin{quote}
\emph{When night came, the girls began to sing.}
\end{quote}

A perfect with no present in classical use: \emph{coepī, coepisse}, participle \emph{coeptus}. It serves as the past of \emph{incipiō}.

\begin{cousinweb}
\textbf{coepī} is the perfect of an old verb \emph{coepiō}, which Lewis \& Short take as \emph{co-} and \emph{apiō}, \textquotedblleft{}to fasten\textquotedblright{}: to join on, to set about.
\end{cousinweb}

\begin{grammarlens}[title={What you've built}]
One more word for telling, explaining and describing at B1.
\end{grammarlens}

\subsection*{Your turn}
\begin{itemize}
\raggedright
  \item \emph{Say it:} \emph{coepī, coepisse}
  \item \emph{Say it:} \emph{coepī, coepisse}, once more
  \item \emph{Say it:} a sentence of your own with \emph{coepī, coepisse}
  \item \emph{Recall:} say the Latin for a moment, then the Latin for to begin, then say \emph{coepī, coepisse} again
\end{itemize}

\begin{tcolorbox}[breakable,colback=teal!4,colframe=teal!35!black,title={Before you move on}]
What does coepī, coepisse mean? (\textquotedblleft{}I began, I have begun\textquotedblright{}.) Say it once more.
\end{tcolorbox}
\section[fīniō, fīnīre]{fīniō, fīnīre — to finish; to set limits}

\label{lesson:LA-C297-finio}



\begin{quote}
Before the new one: say the Latin for to begin, then the Latin for I began.
\end{quote}

\subsection*{You'll want to know: fīniō, fīnīre}
\textbf{fīniō, fīnīre} — \textquotedblleft{}to finish; to set limits\textquotedblright{}.

\begin{quote}
Hōrā nōnā labōrem fīnīvimus.
\end{quote}

\begin{quote}
\emph{We finished the work at the ninth hour.}
\end{quote}

\emph{fīniō, fīnīre, fīnīvī, fīnītum}. From \emph{fīnis} (taught): first 'to mark a boundary', then 'to end'. English \emph{finish}, \emph{finite}.

\begin{cousinweb}
\textbf{fīniō} is built on \emph{fīnis}, \textquotedblleft{}end, boundary\textquotedblright{} (EtymDB): to set an end. English \emph{finish} comes from it through French, and \emph{finite} from its participle.
\end{cousinweb}

\begin{grammarlens}[title={What you've built}]
One more word for telling, explaining and describing at B1.
\end{grammarlens}

\subsection*{Your turn}
\begin{itemize}
\raggedright
  \item \emph{Say it:} \emph{fīniō, fīnīre}
  \item \emph{Say it:} \emph{fīniō, fīnīre}, once more
  \item \emph{Say it:} a sentence of your own with \emph{fīniō, fīnīre}
  \item \emph{Recall:} say the Latin for to begin, then the Latin for I began, then say \emph{fīniō, fīnīre} again
\end{itemize}

\begin{tcolorbox}[breakable,colback=teal!4,colframe=teal!35!black,title={Before you move on}]
What does fīniō, fīnīre mean? (\textquotedblleft{}To finish; to set limits\textquotedblright{}.) Say it once more.
\end{tcolorbox}
\section[perficiō, perficere]{perficiō, perficere — to complete, to carry out}

\label{lesson:LA-C297-perficio}



\begin{quote}
Before the new one: say the Latin for I began, then the Latin for to finish.
\end{quote}

\subsection*{You'll want to know: perficiō, perficere}
\textbf{perficiō, perficere} — \textquotedblleft{}to complete, to carry out\textquotedblright{}.

\begin{quote}
Mūrum novum tandem perfēcērunt.
\end{quote}

\begin{quote}
\emph{At last they completed the new wall.}
\end{quote}

\emph{perficiō, perficere, perfēcī, perfectum}: \emph{per} 'through' + \emph{faciō}. English \emph{perfect}.

\begin{cousinweb}
\textbf{perficiō} is built on \emph{faciō}, \textquotedblleft{}to make, to do\textquotedblright{}, with \emph{per-}, \textquotedblleft{}through\textquotedblright{} (EtymDB): to do right through. English \emph{perfect} comes from its participle.
\end{cousinweb}

\begin{grammarlens}[title={What you've built}]
One more word for telling, explaining and describing at B1.
\end{grammarlens}

\subsection*{Your turn}
\begin{itemize}
\raggedright
  \item \emph{Say it:} \emph{perficiō, perficere}
  \item \emph{Say it:} \emph{perficiō, perficere}, once more
  \item \emph{Say it:} a sentence of your own with \emph{perficiō, perficere}
  \item \emph{Recall:} say the Latin for I began, then the Latin for to finish, then say \emph{perficiō, perficere} again
\end{itemize}

\begin{tcolorbox}[breakable,colback=teal!4,colframe=teal!35!black,title={Before you move on}]
What does perficiō, perficere mean? (\textquotedblleft{}To complete, to carry out\textquotedblright{}.) Say it once more.
\end{tcolorbox}
\section[intermittō, intermittere]{intermittō, intermittere — to break off, to interrupt}

\label{lesson:LA-C297-intermitto}



\begin{quote}
Before the new one: say the Latin for to finish, then the Latin for to complete.
\end{quote}

\subsection*{You'll want to know: intermittō, intermittere}
\textbf{intermittō, intermittere} — \textquotedblleft{}to break off, to interrupt\textquotedblright{}.

\begin{quote}
Propter pluviam labōrem intermīsimus.
\end{quote}

\begin{quote}
\emph{Because of the rain we broke off the work.}
\end{quote}

\emph{intermittō, intermittere, intermīsī, intermissum}: to put a gap into something. Caesar uses it of a battle or a march that is broken off.

\begin{cousinweb}
\textbf{intermittō} is built on \emph{mittō}, \textquotedblleft{}to send, to let go\textquotedblright{}, with \emph{inter}, \textquotedblleft{}between\textquotedblright{} (Lewis \& Short): to leave a gap. English \emph{intermission} and \emph{intermittent} come from it.
\end{cousinweb}

\begin{grammarlens}[title={What you've built}]
One more word for telling, explaining and describing at B1.
\end{grammarlens}

\subsection*{Your turn}
\begin{itemize}
\raggedright
  \item \emph{Say it:} \emph{intermittō, intermittere}
  \item \emph{Say it:} \emph{intermittō, intermittere}, once more
  \item \emph{Say it:} a sentence of your own with \emph{intermittō, intermittere}
  \item \emph{Recall:} say the Latin for to finish, then the Latin for to complete, then say \emph{intermittō, intermittere} again
\end{itemize}

\begin{tcolorbox}[breakable,colback=teal!4,colframe=teal!35!black,title={Before you move on}]
What does intermittō, intermittere mean? (\textquotedblleft{}To break off, to interrupt\textquotedblright{}.) Say it once more.
\end{tcolorbox}
